\chapter{Conformal fields, states, and operator products}
\label{ch:cft}

Conformal field theory organizes infinitely many local operators using
symmetry. In two dimensions the organization is unusually powerful:
local conformal symmetry separates into holomorphic and antiholomorphic
parts, whose generators form two commuting Virasoro algebras. On the
Euclidean plane $\bar z$ is the complex conjugate of $z$; treating the
two parts separately is a way to organize the symmetry. This chapter
develops the pieces of that structure needed
later. The aim is to read an operator product expansion (OPE), compute
a descendant norm,
translate a conformal weight into a finite-circle energy, and recognize
which information a character does and does not contain.

A \emph{local field} $\phi(z,\bar z)$ is an operator insertion at a
point. Its measurable content is encoded in correlation functions such
as $\langle\phi_1(z_1,\bar z_1)\phi_2(z_2,\bar z_2)\rangle$, the
vacuum expectation value of the indicated insertions. Fields are
generally operator-valued distributions, so coincident products need a
prescription; separated-point correlators and the OPE provide one.
Conformal transformations preserve angles locally. In two dimensions
an orientation-preserving such transformation is locally a holomorphic
change of coordinate with nonzero derivative. We first describe the
physics of its action, and leave the formal algebraic definition of
fields and their modules to Chapter~\ref{ch:vertex_algebras}.

\section{From scale transformations to radial quantization}
\index{radial quantization}\index{state--operator correspondence}\index{conformal spin}

Write a Euclidean point as $z=x+iy$. A rotation and a dilation take
$z$ to $z'=\lambda e^{i\theta}z$, with $\lambda>0$. For a field
of definite scaling dimension and spin, our coordinate-change
convention is
\[
 \phi'(z',\bar z')=\lambda^{-\Delta}e^{-i s_{\rm spin}\theta}
                         \phi(z,\bar z).
\]
The \emph{holomorphic and antiholomorphic conformal weights}
$(h,\bar h)$ are the two numbers satisfying
\begin{equation}
 \Delta=h+\bar h,\qquad s_{\rm spin}=h-\bar h.
 \label{eq:cft-dimension-spin}
\end{equation}
The bar on $\bar h$ labels the antiholomorphic weight; it does not
mean complex conjugation of $h$.
For a bosonic local field, a full rotation must return the same field,
so $s_{\rm spin}\in\Z$. A scalar has $s_{\rm spin}=0$, a stronger
condition than integer spin. A \emph{chiral} field depends only on
$z$ in separated-point correlators; an \emph{antichiral} field depends
only on $\bar z$. A holomorphic field of weights $(1,0)$ is called
a chiral current. Its contour integral generates a symmetry. A field
of weights $(1,1)$ is instead a dimension-two scalar.
These examples will prevent several otherwise easy mistakes when
reading numerical spectra.
The integer-spin requirement applies to full bosonic local fields.
A chiral factor used to build one can have noninteger weight and
need not itself define a single-valued local field of the full theory.

Radial quantization treats circles about the origin as equal-time
slices. With an arbitrary reference length set to one, radial time is
$t=\log|z|$; moving outwards is moving forwards in this time. The
path integral over a disk sums field configurations in its interior
with fixed boundary values; as a function of those boundary values
it prepares a state on the boundary circle. An empty disk prepares the
vacuum $|0\rangle$. A field inserted at the origin creates a state,
\[
 |\phi\rangle=\lim_{z,\bar z\to0}\phi(z,\bar z)|0\rangle.
\]
Here the ket is a vector in the state space, whereas $\phi(z,\bar z)$
is a field acting on states. The limit means the state prepared by the
insertion; it does not assert that the field is a bounded operator.
The \emph{state--operator correspondence} identifies local fields
with states in this way. We use it for theories with a normalizable
vacuum and a discrete spectrum with finite degeneracies at fixed
energy. More general theories require qualifications.

Write $L_0$ and $\bar L_0$ for the operators measuring the two
weights:
\[
 L_0|\phi\rangle=h|\phi\rangle,\qquad
 \bar L_0|\phi\rangle=\bar h|\phi\rangle.
\]
The generator of radial dilations is $L_0+\bar L_0$; rotations are
generated by $L_0-\bar L_0$. The next section constructs these
operators as members of two infinite families. Expanding a product of
insertions inside a disk into states on its boundary will then give
the OPE.

\section{The stress tensor and its central term}
\index{stress tensor}\index{Virasoro algebra}\index{central charge}

The holomorphic stress tensor $T(z)$ generates infinitesimal conformal
transformations. It has weight two. An \emph{operator product expansion}
expresses two nearby insertions as a sum of single insertions with
coordinate-dependent coefficients. Such an equality is read inside
correlators, with the other insertions farther away. The stress
tensor's self-OPE is
\begin{equation}
 T(z)T(w)\sim \frac{c/2}{(z-w)^4}
       +\frac{2T(w)}{(z-w)^2}+\frac{\partial T(w)}{z-w}.
 \label{eq:cft-TT}
\end{equation}
The first term means $(c/2)(z-w)^{-4}$ times the identity field
$\one$, whose insertion changes nothing. The derivative
$\partial T(w)$ means $\partial_w T(w)$.
The symbol $\sim$ retains the singular part. Regular terms exist and
matter for composite fields, but they do not enter the contour
commutator that we are about to compute. The number $c$ is the chiral
central charge. It measures the central extension of conformal
symmetry and controls the vacuum's response to a change of geometry.

Expand the field in a Laurent series. Its coefficient operators,
called \emph{modes}, are
\[
 T(z)=\sum_{n\in\Z}L_nz^{-n-2},\qquad
 L_n=\frac1{2\pi i}\oint z^{n+1}T(z)\,dz.
\]
Each $L_n$ is a linear operator on states. The contour is a small
counterclockwise circle about the origin; the integral picks out the
coefficient of $z^{-1}$. Radially ordered products place the operator
at larger radius to the left. Subtracting the two radial orderings
deforms the contours into one about the collision $z=w$.
Evaluating that contour gives
\begin{equation}
 [L_m,L_n]=(m-n)L_{m+n}
       +\frac{c}{12}(m^3-m)\delta_{m+n,0}.
 \label{eq:cft-virasoro}
\end{equation}
Here $[A,B]=AB-BA$ and the Kronecker symbol $\delta_{m+n,0}$ is
one if $m+n=0$ and zero otherwise. The scalar central term multiplies
the identity operator on the state space; it commutes with all modes.
These commutation relations define the \emph{Virasoro algebra}.
For example, the fourth-order pole contributes the third derivative
of $z^{m+1}$ divided by $3!$. Its residue against $w^{n+1}$ is
$c(m+1)m(m-1)\delta_{m+n,0}/12$, explaining both the cubic polynomial
and its vanishing for $m=0,\pm1$. These three modes generate the
ordinary global conformal algebra without a central term.
It is the Lie algebra $\mathfrak{sl}_2$, generated here by
$L_{-1},L_0,L_1$; geometrically these generate translations,
dilations, and special conformal transformations. In particular,
$[L_0,L_n]=-nL_n$, so $L_n$ shifts a weight by $-n$.

There is a commuting antiholomorphic copy with central charge $\bar c$.
For the nonchiral theories studied in this book, $c=\bar c$.
When we say that a theory has central charge eight, we mean
$(c,\bar c)=(8,8)$, not that their sum is eight.

\begin{definition}
A \emph{Virasoro primary}, or \emph{highest-weight state}, of weight
$h$ satisfies
$L_0|\phi\rangle=h|\phi\rangle$ and
$L_n|\phi\rangle=0$ for every $n>0$. Its field obeys
\begin{equation}
 T(z)\phi(w)\sim\frac{h\phi(w)}{(z-w)^2}
                         +\frac{\partial\phi(w)}{z-w}.
 \label{eq:cft-Tprimary}
\end{equation}
A quasiprimary state is annihilated by $L_1$ and has definite weight;
it need not be annihilated by $L_2,L_3,\ldots$.
\end{definition}

Thus a primary has no lower-weight state obtained by applying a
positive Virasoro mode. Equivalently its field transforms under a
local coordinate change as
$\phi'(f(z))=(f'(z))^{-h}\phi(z)$, with a local branch chosen for
the power. A full primary obeys the corresponding conditions in both
chiralities. Unless a larger algebra is specified, ``primary'' in
this chapter always means Virasoro primary. The stress tensor is
quasiprimary, but for $c\ne0$ its fourth-order self-OPE pole prevents
it from being primary.

Contour integration of \eqref{eq:cft-Tprimary} gives
\[
 [L_m,\phi(z)]
 =z^{m+1}\partial\phi(z)+(m+1)h z^m\phi(z).
\]
In particular, $L_{-1}$ differentiates a field. Acting with negative
Virasoro modes generates \emph{descendants}, the symmetry-related
states above a primary. The descendant
$L_{-n_1}\cdots L_{-n_r}|\phi\rangle$ has weight
$h+n_1+\cdots+n_r$, for positive integers $n_i$. The sum
$n_1+\cdots+n_r$ is its \emph{level}, or grade. A representation,
also called a module, of this algebra is a vector space on which
the $L_n$ act and satisfy
\eqref{eq:cft-virasoro}. The \emph{Verma module} is the universal
highest-weight representation: start with one $|h\rangle$ and impose
only the Virasoro relations and highest-weight conditions. Ordered
words with $n_1\ge\cdots\ge n_r\ge1$ give its basis. A physical
representation can be a quotient of this space by additional
relations. Null states are crucial examples of such relations.

\section{Unitarity and descendant norms}
\index{null state}\index{Gram matrix!Virasoro}

Unitarity means that the physical state space has a positive inner
product compatible with time evolution and radial adjoints. For the
Virasoro modes this compatibility reads
$L_n^\dagger=L_{-n}$, where the dagger denotes the adjoint with
respect to that inner product. Radial reflection exchanges the inside
and outside of the unit circle. For a chiral primary of real weight
$h$, it represents the bra by an insertion at infinity:
\[
 \langle\phi|=\lim_{R\to\infty}R^{2h}
                    \langle0|\phi^\dagger(R).
\]
The conjugate chiral field $\phi^\dagger$ is still a field in the
holomorphic theory; it is not the independent antiholomorphic field.
We abbreviate this limit as $\phi^\dagger(\infty)$. Unit normalization
means $\langle\phi^\dagger(\infty)\phi(0)\rangle=1$.
For a full primary on the positive real axis, the factor is
$R^{2(h+\bar h)}$.

If $|h\rangle$ is a unit primary, then
\[
 \norm{L_{-1}|h\rangle}^2
 =\langle h|[L_1,L_{-1}]|h\rangle=2h.
\]
Consequently $h\ge0$. At $h=0$ the translation descendant is null.
In a positive semidefinite representation, a \emph{null state} is
a vector of norm zero and is orthogonal to every state. Such vectors
are set equal to zero to obtain the physical Hilbert space. Before
this quotient they can be nonzero formal vectors. In an indefinite
inner product, zero norm alone does not imply orthogonality to every
state; there one must specify the radical of the form, the subspace
orthogonal to every vector, when taking a quotient.
The vacuum obeys the stronger geometric condition
$L_n|0\rangle=0$ for $n\ge-1$.

\begin{example}[The level-two Gram matrix]
\label{ex:cft-level2}
At level two the Verma module has the ordered basis
$v_1=L_{-2}|h\rangle,v_2=L_{-1}^2|h\rangle$.
The \emph{Gram matrix} records their pairwise inner products,
$(G_2)_{ij}=\langle v_i|v_j\rangle$. It is
\begin{equation}
 G_2(h,c)=
 \begin{pmatrix}
 4h+c/2&6h\\
 6h&4h(2h+1)
 \end{pmatrix}.
 \label{eq:cft-gram2}
\end{equation}
The first entry is $\langle h|[L_2,L_{-2}]|h\rangle$.
For the off-diagonal entry, use
$[L_2,L_{-1}]=3L_1$ and
$L_1L_{-1}|h\rangle=2h|h\rangle$ to obtain $6h$.
Finally,
\[
 L_1L_{-1}^2|h\rangle=2(2h+1)L_{-1}|h\rangle,
\]
which gives the bottom-right entry. Its determinant is
\[
 \det G_2=2h\bigl(16h^2+2(c-5)h+c\bigr).
\]
For a coefficient column $u$, the corresponding state has squared
norm $u^\dagger G_2u$. A nonzero vector in the kernel of $G_2$
therefore describes a linear combination orthogonal to this whole
level. Distinct levels are orthogonal because $L_0$ is self-adjoint,
so a zero determinant detects a null relation. Positivity of every such Gram matrix
at every level is much stronger than positivity of the primary norm.
\end{example}

For the vacuum, $L_{-1}|0\rangle=0$, but
\begin{equation}
 \norm{L_{-2}|0\rangle}^2=c/2.
 \label{eq:cft-stress-norm}
\end{equation}
The stress state is therefore $|T\rangle=L_{-2}|0\rangle$, and its
unit-normalized version, when $c>0$, is $\sqrt{2/c}\,|T\rangle$.
This simple norm
will later be the basis of a central-charge estimator for the chain.

Null vectors must be quotiented out to obtain the physical state space.
Counting all words in the negative modes without doing this would
overcount states. Conversely, an observed deficit in a character is
not automatically a new null relation. A character is a graded state
count, defined in Section~\ref{sec:cft-characters}; one first has to specify
which highest-weight module is being counted.

\section{What an operator product specifies}
\index{operator product expansion}

\subsection{Choose the chiral algebra before naming its modules}
\label{sec:cft-chiral-modules}
\index{chiral algebra}\index{module!chiral}

The stress tensor need not be the only holomorphic field whose modes
organize the states. A \emph{chiral algebra} is a specified collection
of mutually local holomorphic fields containing the identity and
stress tensor and closed under derivatives and OPE products. Here
mutual locality means that exchanging these algebra fields produces
no nontrivial branch monodromy: their OPEs use integer powers, with
the usual bosonic locality relations. Chapter~\ref{ch:vertex_algebras} makes
this precise using vertex operator algebras. The chiral algebra
generated by $T$, whose modes form the Virasoro Lie algebra, is the
minimal example. Adjoining currents or other
holomorphic fields gives an \emph{extended chiral algebra}.

A \emph{module}, or representation, for a chosen chiral algebra is
a vector space $M$ on which all the modes of its fields act, obeying
the same OPE identities as in the algebra. For example, the modes
$L_n:M\to M$ must satisfy the Virasoro commutators. If a current
$J$ belongs to the algebra, its modes must also act on $M$ and
satisfy their OPE relations with $T$, with each other, and with the
other fields. Thus $M$ consists of states together with specified
operator actions. It is neither a single field nor merely a number
of states.

A module is \emph{irreducible}, also called \emph{simple}, if it
has no nonzero proper subspace preserved by every algebra mode.
An irreducible representation type is often called a \emph{sector}.
The \emph{vacuum module} contains $|0\rangle$ and the states
obtained from it using the algebra fields. Other modules need not
contain a vacuum. A Virasoro highest-weight module is the familiar
example generated by the negative $L_n$ from a primary, after any
required quotient. A module for a larger algebra can contain many
Virasoro primaries: its extra modes connect their Virasoro families.
Consequently the number of irreducible modules depends on which
chiral algebra we choose.

It is also useful to distinguish an algebra field from a general
chiral field. The former acts within each module. A charged chiral
field can instead carry a state from one module into another, and
its products may involve fractional powers. Such fields are described
by \emph{intertwining operators}: field-valued maps between modules
compatible with the algebra action. We will use their basic OPE
meaning here; their precise definition comes in
Chapter~\ref{ch:vertex_algebras}. A full local field combines holomorphic and
antiholomorphic data so that its correlators have the required
single-valuedness.

\subsection{Coefficients, descendants, and fusion channels}

If two local fields are close compared with every other separation,
their product can be expanded in local fields at a common point.
For chiral primaries the schematic form is
\begin{equation}
 \phi_i(z)\phi_j(0)
 =\sum_k C_{ij}^{\ k}z^{h_k-h_i-h_j}
       \bigl(\phi_k(0)+\hbox{descendants}\bigr).
 \label{eq:cft-primary-ope}
\end{equation}
The descendant coefficients are constrained by conformal symmetry;
the primary coefficients $C_{ij}^{\ k}$ contain dynamical information.
The power of $z$ balances the conformal weights. Each level-$n$
descendant inside the parentheses comes with an additional $z^n$.
For instance, if the exchanged primary has weight $h_k>0$, the
level-one term is
\[
 \phi_k(0)+z\,\frac{h_k+h_i-h_j}{2h_k}
                    \partial\phi_k(0)+\cdots.
\]
Its coefficient follows by applying $L_1$ and using
$L_1L_{-1}|h_k\rangle=2h_k|h_k\rangle$. At higher levels one
solves analogous linear equations using descendant Gram matrices,
with null states removed. This is why symmetry fixes descendant
coefficients once the primary coupling is chosen.

An \emph{OPE channel} specifies the representation exchanged between
the two insertions and the rest of the correlator. For a chosen chiral
algebra, the allowed representation types form its \emph{fusion
rules}. A fusion coefficient counts independent compatible chiral
couplings, just as a junction multiplicity did in
Chapter~\ref{ch:fusion}; it does not count descendants. If more
than one such coupling reaches the same output representation, one
must attach an extra intertwiner label and choose a basis of that
finite-dimensional coupling space. Equation~\eqref{eq:cft-primary-ope}
suppresses these labels and resolves fields into Virasoro families.

Choose an orthonormal primary basis, so that
$\langle\phi_k^\dagger(\infty)\phi_l(0)\rangle=\delta_{kl}$.
Then the three-point coefficient can be read as
\[
 C_{ij}^{\ k}
 =\langle\phi_k^\dagger(\infty)\phi_i(1)\phi_j(0)\rangle.
\]
Conformal invariance then fixes all coordinate dependence:
\begin{equation}
 \langle\phi_i(z_1)\phi_j(z_2)\phi_k(z_3)\rangle
 =\frac{C_{ijk}}
 {z_{12}^{h_i+h_j-h_k}
  z_{13}^{h_i+h_k-h_j}
  z_{23}^{h_j+h_k-h_i}}.
 \label{eq:cft-threepoint}
\end{equation}
Here $z_{ij}=z_i-z_j$. For full fields there is an analogous
antiholomorphic factor.
The all-lower coefficient $C_{ijk}$ uses a two-point pairing with
the third field: if $g_{lk}=\langle\phi_l(\infty)\phi_k(0)\rangle$,
then, with consistent field ordering and branches,
$C_{ijk}=\sum_l C_{ij}^{\ l}g_{lk}$.
For charged fields this pairing typically pairs opposite charges;
lowering the index is not the same as ignoring a dagger.
Locality and branch conventions must be respected when the chiral
powers are not integers. A chiral contribution can be multivalued even
when the full correlation function is single-valued. We define the
individual contributions called conformal blocks in
Section~\ref{sec:cft-blocks}.

Field phases change individual complex coefficients. Under
$\phi_i\mapsto e^{i\theta_i}\phi_i$, the coefficient
$C_{ij}^{\ k}$ changes by
$e^{i(\theta_i+\theta_j-\theta_k)}$. Norms do not fix these phases.
Products of coefficients can cancel every phase and hence contain
physical information. A real Hamiltonian or a positive Hilbert-space
metric does not require every mixed coefficient in every charged
basis to be real. The overlap chapter will use explicit invariants
to separate phase conventions from measured relative phases.

\begin{example}[The universal stress coefficient]
Let $\phi$ be a unit chiral primary of weight $h$ and let $c>0$.
The contribution of the \emph{Virasoro vacuum family}, consisting of
the identity and its Virasoro descendants, to
$\phi(z)\phi^\dagger(0)$ begins
\begin{equation}
 \left.\phi(z)\phi^\dagger(0)\right|_{\text{Virasoro vacuum}}
 =z^{-2h}\left(\one+\frac{2h}{c}z^2T(0)+\cdots\right).
 \label{eq:cft-stress-coefficient}
\end{equation}
To derive the coefficient, insert $T(w)$ far from the pair.
A \emph{Ward identity} expresses an insertion of a symmetry current
through its action on the other fields. Summing the poles of
\eqref{eq:cft-Tprimary} gives the stress-tensor Ward identity on the
plane,
\[
 \left\langle T(w)\prod_i\phi_i(z_i)\right\rangle
 =\sum_i\left(\frac{h_i}{(w-z_i)^2}
       +\frac{1}{w-z_i}\partial_{z_i}\right)
       \left\langle\prod_i\phi_i(z_i)\right\rangle.
\]
Apply it to the two-point function $(z_1-z_2)^{-2h}$, and then set
$z_1=z,z_2=0$. The result is
$h z^{2-2h}/[w^2(w-z)^2]$, whose leading large-$w$ term is
$h z^{2-2h}w^{-4}$.
The proposed OPE instead gives
$A z^{2-2h}\langle T(w)T(0)\rangle
=A(c/2)z^{2-2h}w^{-4}$, so $A=2h/c$.
The coupling to the unit stress state is therefore
$h\sqrt{2/c}$, whose square is $2h^2/c$.
For $h=2,c=8$, this coefficient is one.
Other families may also appear in the OPE. In particular a current
can belong to the vacuum module of an extended algebra while being
a separate Virasoro primary. Such a current contribution has not
been included in the displayed Virasoro-vacuum projection.
\end{example}

\section{Free bosons as a complete first laboratory}
\index{free boson}\index{compact boson}

A chiral free boson can be normalized by
$X(z)X(w)\sim-\log(z-w)$. This formula specifies its Gaussian
two-point contraction, the part used when evaluating products by
pairwise contractions. Because of the logarithm, $X$ itself is an
auxiliary chiral field with branch and zero-mode subtleties; its
derivative is the single-valued current we use in the algebra.
Set $J=i\partial X$. Differentiating
the logarithm with the appropriate factors of $i$ gives
\[
 J(z)J(w)\sim\frac1{(z-w)^2}.
\]
\emph{Normal ordering}, denoted by colons, subtracts the singular
self-contraction before taking coincident points. For example,
\[
 :JJ:(w)=\lim_{z\to w}
       \left(J(z)J(w)-\frac{\one}{(z-w)^2}\right).
\]
The limit is understood inside separated-point correlators.
For these Gaussian fields, \emph{Wick's rule} expresses a product
of normal-ordered fields as the sum over all pairwise contractions
between them. The stress tensor
$T=\tfrac12:JJ:$ has central charge one: in $T(z)T(w)$ there are
two double contractions, giving
$\tfrac14\cdot2(z-w)^{-4}=\tfrac12(z-w)^{-4}$.
The single contractions give the required $2T$ and $\partial T$
terms. Building a stress tensor quadratically from currents is called
a \emph{Sugawara construction}; here it is just this elementary
free-boson calculation.

For a real charge $p$, the normal-ordered exponential
$V_p=:e^{ipX}:$ is obtained by expanding the exponential and
normal-ordering each term. Wick's rule resums the contractions
between two exponentials into
$\exp[-pq\langle X(z)X(w)\rangle]=(z-w)^{pq}$.
It gives
\begin{align}
 J(z)V_p(w)&\sim\frac{p}{z-w}V_p(w),\\
 V_p(z)V_q(w)&=(z-w)^{pq}
       \bigl(V_{p+q}(w)+\hbox{oscillator descendants}\bigr),\\
 h(V_p)&=p^2/2.
 \label{eq:cft-boson-weights}
\end{align}
The exponent in the OPE agrees with
$h(p+q)-h(p)-h(q)=pq$. This agreement is a useful elementary check
on every later lattice-exponential formula.

The weight follows from the double contraction in
$T(z)V_p(w)$: its double pole is $(p^2/2)V_p(w)$.
The single contractions give $\partial V_p(w)$ in the simple pole,
so $V_p$ is a Virasoro primary.

\subsection{Current modes and a concrete chiral module}

Expand the current as
\[
 J(z)=\sum_{n\in\Z}J_nz^{-n-1},\qquad
 J_n=\frac{1}{2\pi i}\oint z^nJ(z)\,dz.
\]
The $JJ$ OPE gives $[J_m,J_n]=m\delta_{m+n,0}\,\mathrm{id}$.
Wick's rule gives
$T(z)J(w)\sim J(w)/(z-w)^2+\partial J(w)/(z-w)$,
so $J$ is a weight-one primary and
$[L_0,J_n]=-nJ_n$. Thus $J_{-r}$ raises the weight by $r$ for
$r>0$; it is a creation operator. The operator $J_0$, called the
\emph{zero mode} of the current, measures charge rather than
raising or lowering weight.

Let $|p\rangle=V_p(0)|0\rangle$. Its defining conditions are
\[
 J_0|p\rangle=p|p\rangle,\qquad
 J_n|p\rangle=0\ (n>0),\qquad
 L_0|p\rangle=\frac{p^2}{2}|p\rangle.
\]
The \emph{Fock module} $F_p$ is the vector space with basis
\[
 \prod_{r\ge1}J_{-r}^{k_r}|p\rangle,
 \qquad k_r\in\Z_{\ge0},\quad\text{only finitely many }k_r\ne0.
\]
The negative modes act by adding factors, positive modes act through
the commutators, and $J_0$ acts by $p$ throughout. For example,
$J_1J_{-1}|p\rangle=|p\rangle$ and
$J_1J_{-2}|p\rangle=0$. This is a concrete module for the chiral
algebra generated by $J$ and $T=\tfrac12:JJ:$.
Its basis state above has weight $p^2/2+\sum_r r k_r$.
The term \emph{oscillator descendants} in the exponential OPE means
precisely states formed by these current creation modes.

This example also distinguishes two kinds of descendant. The
Sugawara formula gives
$L_{-1}|p\rangle=pJ_{-1}|p\rangle$. At $p=0$ the Virasoro
translation descendant vanishes, but $J_{-1}|0\rangle$ does not.
That state is the current itself, a new Virasoro primary inside the
current algebra's vacuum module. We will count the states of $F_p$
explicitly in Section~\ref{sec:cft-characters}.

\subsection{Charge conservation and compactification}

The current Ward identity is
\[
 \left\langle J(z)\prod_jV_{p_j}(z_j)\right\rangle
 =\sum_j\frac{p_j}{z-z_j}
             \left\langle\prod_jV_{p_j}(z_j)\right\rangle.
\]
Integrating around all insertions computes their total $J_0$ charge.
Since the vacuum at infinity has charge zero, a nonzero vacuum
correlator requires $\sum_jp_j=0$. For charges summing to zero,
\begin{equation}
 \left\langle\prod_jV_{p_j}(z_j)\right\rangle
 =\prod_{i<j}(z_i-z_j)^{p_ip_j};
 \qquad \sum_jp_j=0.
 \label{eq:cft-gaussian}
\end{equation}
Otherwise the correlator vanishes. With several bosons, replace
$p_ip_j$ by the Euclidean inner product of charge vectors. A lattice
theory also needs \emph{cocycles}: consistent phase factors in
multiplying its charged fields. They fix exchange signs that the
bare exponentials alone do not supply; the next chapter explains
their role.

\begin{example}[A charged four-point function]
For $h=p^2/2$, place charges $(-p,p,-p,p)$ at
$(\infty,1,x,0)$. Equation~\eqref{eq:cft-gaussian} gives
\[
 \langle V_{-p}(\infty)V_p(1)V_{-p}(x)V_p(0)\rangle
 =x^{-p^2}(1-x)^{-p^2}.
\]
Near $x=0$, expand $(1-x)^{-p^2}=1+p^2x+\cdots$.
The coefficient $p^2$ is the square of the coupling to the unit
current $J$, as follows independently from its Ward identity.
Indeed the full current-algebra OPE starts
\[
 V_p(z)V_{-p}(0)
 =z^{-p^2}\bigl(\one+pzJ(0)+\cdots\bigr).
\]
The current contribution explains the linear term in the four-point
function; the Virasoro vacuum family alone has no level-one state.
This is a small version of the charge-expansion calculation used for
the defining Haagerup field $u$.
\end{example}

A \emph{compact} full boson has a real field identified by
$X_{\rm full}\equiv X_{\rm full}+2\pi R$, with $R>0$ its target
circle's radius. This target circle differs from the spatial circle
on which the quantum theory is placed. Its zero-mode wavefunctions
have integer momentum label $m$, and a field configuration can wind
$n$ times around the target circle while going once around space.
The pairs $(m,n)\in\Z^2$ form the momentum--winding lattice.
In the radius convention used here, their left and right charges are
\[
 p_L=\frac mR+\frac{nR}{2},\qquad
 p_R=\frac mR-\frac{nR}{2},\qquad m,n\in\Z,
\]
with $h=p_L^2/2$, $\bar h=p_R^2/2$ for the ground state in each
charge sector. Oscillator excitations add nonnegative integer levels
to each weight. Consequently the ground-state dimension and spin are
\[
 \Delta=\frac{m^2}{R^2}+\frac{n^2R^2}{4},\qquad
 s_{\rm spin}=mn\in\Z.
\]
The radius changes dimensions while preserving the central charge.
There are infinitely many charge sectors for the algebra generated
only by the current. When $R^2$ is rational in this convention, an
extension by suitable charged holomorphic fields can organize them
into finitely many irreducible modules. Thus ``rational radius'' here
refers to $R^2$, not $R$; the self-dual example $R=\sqrt2$ qualifies.
Rationality is a property of a specified chiral algebra and its
representations, not a synonym for having some rational-looking
low dimensions.

\section{From the plane to the cylinder}
\index{cylinder map}\index{Casimir energy}

Set $z=e^{2\pi w/L}$, where $w=vt_E+ix$, $t_E$ is Euclidean
time, and $x$ has spatial period $L$. The constant $v>0$ converts
time to length. Advancing $x$ by $L$ leaves $z$ unchanged; advancing
$t_E$ moves radially outwards. This is why the map converts radial
evolution on the plane into time evolution on a cylinder.
Unlike a primary, the stress tensor acquires a geometry-dependent
term under a coordinate change. That term involves the following
combination of derivatives, called the \emph{Schwarzian derivative}:
\[
 T_{\rm cyl}(w)=\left(\frac{dz}{dw}\right)^2T(z)
                   +\frac c{12}\{z,w\},
 \qquad
 \{z,w\}=\frac{z'''}{z'}-\frac32\left(\frac{z''}{z'}\right)^2.
\]
Primes mean derivatives with respect to $w$.
For this exponential map,
$\{z,w\}=-\tfrac12(2\pi/L)^2$. The cylinder Hamiltonian and
momentum are therefore
\begin{align}
 H_{\rm cyl}&=\frac{2\pi v}{L}
       \left(L_0+\bar L_0-\frac{c+\bar c}{24}\right),
 \label{eq:cft-cylinder}\\
 P_{\rm cyl}&=\frac{2\pi}{L}
       \left(L_0-\bar L_0-\frac{c-\bar c}{24}\right).
\end{align}
For equal chiral central charges and a vacuum of weight zero,
$E_0=-\pi vc/(6L)$. A lattice Hamiltonian adds an extensive
nonuniversal energy $e_\infty L$, where $e_\infty$ is the bulk
ground-state energy per unit length. The negative $1/L$ vacuum term
is the \emph{Casimir energy}, the finite-size shift created by the
change of geometry. Terms that vanish faster at large $L$ can come
from irrelevant perturbations, defined in
Section~\ref{sec:cft-perturbations}. The lattice
excitation gaps approach $2\pi v\Delta/L$ only when a CFT limit
and the appropriate scaling regime have been reached.

For a state at left and right descendant levels $N,\bar N$ above
weights $(h,\bar h)$, replace $\Delta$ in that gap by
$h+\bar h+N+\bar N$. For $c=\bar c$, its momentum relative to
the vacuum is $2\pi(h-\bar h+N-\bar N)/L$. Thus characters will
count precisely the excitations whose circle energies these formulas
describe. We use units with $\hbar=1$.

The velocity $v$ is a conversion between lattice time and continuum
space; it is not fixed by conformal symmetry. An energy fit that
extracts the coefficient of $1/L$ in the vacuum energy measures
$vc$, not $c$ separately. A second observable is needed to calibrate
$v$. The distinction becomes crucial when comparing a theory of
central charge exactly two with one near $2.025$.

\section{Characters and the torus}
\label{sec:cft-characters}
\index{character}\index{modular invariance}

\subsection{What a character counts}

Recall from Section~\ref{sec:cft-chiral-modules} that a chiral module
$M$ is a vector space carrying the modes of a specified chiral
algebra. To define the character used here, suppose that $L_0$ is
diagonalizable on $M$, has a lowest eigenvalue $h$, and gives a
decomposition
\[
 M=\bigoplus_{n\ge0}M_{h+n},\qquad
 M_{h+n}=\{v\in M:L_0v=(h+n)v\},\qquad
 a_n=\dim M_{h+n}<\infty.
\]
Some of these eigenspaces may be zero. This is the ordinary,
positive-energy, integer-spaced grading appropriate to the untwisted
modules in our examples. ``Untwisted'' means that the algebra fields
return to themselves after going around the insertion, without an
extra symmetry transformation. Twisted sectors can have different
mode spacings. The integer $n$ is the grade, and $M_h$ is the
\emph{top space}, the lowest-weight eigenspace. Even an irreducible
module can have $a_0>1$.

The \emph{character} is the generating function for these graded
dimensions, with the cylinder vacuum-energy shift included:
\begin{equation}
 \chi_M(\tau)=\Tr_M q^{L_0-c/24},\qquad
 q=e^{2\pi i\tau},\qquad \Im\tau>0.
 \label{eq:cft-character}
\end{equation}
Here $\Tr_M$ is the sum of ordinary linear-algebra traces over the
finite-dimensional graded pieces of $M$. On each of the $a_n$
basis states at weight $h+n$, the traced
operator acts by $q^{h+n-c/24}$. Consequently
\[
 \chi_M(\tau)=q^{h-c/24}\sum_{n\ge0}a_nq^n.
\]
The variable $q$ records weight; it is not the field coordinate $z$.
The condition $\Im\tau>0$ makes $|q|<1$. We may first regard the
sum as a formal series. Analytic convergence also requires control
of the growth of $a_n$; finite-dimensional grades alone do not prove
it. The free-boson example below converges, as do the rational models
used later under their stated hypotheses.

The integer $a_0$ counts lowest-weight states. A quantum dimension,
in the sense introduced in Chapter~\ref{ch:fusion}, instead describes
fusion and the value of a closed categorical line. When a category of
modules is available these are still different notions of dimension.
Likewise a character records multiplicities, but does not specify the
matrices of every mode action or the OPE coefficients.

\begin{example}[Counting a free-boson module]
\label{ex:cft-fock-character}
For $F_p$ the lowest weight is $h=p^2/2$ and $c=1$.
The first three grades have the following bases:
\[
\begin{array}{c|l|c}
 n&\text{states}&a_n\\ \hline
 0&|p\rangle&1\\
 1&J_{-1}|p\rangle&1\\
 2&J_{-2}|p\rangle,\ J_{-1}^2|p\rangle&2.
\end{array}
\]
At grade three the occupation patterns are $3$, $2+1$, and
$1+1+1$, giving three states. More generally, $k_r$ copies of
$J_{-r}$ contribute grade $r k_r$. Summing each occupation number
independently gives
\begin{align*}
 \chi_{F_p}(\tau)
 &=q^{p^2/2-1/24}\prod_{r\ge1}
                 (1+q^r+q^{2r}+\cdots)\\
 &=q^{p^2/2-1/24}\prod_{r\ge1}(1-q^r)^{-1}\\
 &=q^{p^2/2-1/24}(1+q+2q^2+3q^3+5q^4+\cdots).
\end{align*}
The grade-four count comes from the five partitions
$4$, $3+1$, $2+2$, $2+1+1$, and $1+1+1+1$.
The product converges for $|q|<1$ since the sum of the absolute
values $\sum_{r\ge1}|q|^r$ is finite.

The modules $F_p$ and $F_{-p}$ have the same character, although
for $p\ne0$ the current zero mode has different eigenvalues on them.
Thus even a complete character need not determine the module.
At $p=0$ the $q$ term counts the current state; it is present in
the current algebra's vacuum module even though the Virasoro vacuum
has no level-one descendant.
\end{example}

\subsection{Pairing chiral sectors on a torus}

A torus of complex structure $\tau$ is
$\C/(\Z+\tau\Z)$: identify points differing by an integer
combination of the two periods $1$ and $\tau$. Taking
$\Im\tau>0$ fixes the orientation. It can be viewed as a spatial
circle evolved through Euclidean time and glued back to itself,
possibly after a spatial translation.

In the rational setting, finitely many irreducible chiral modules
are available and the relevant state spaces decompose into direct
sums of them. Complete reducibility, rather than finiteness alone,
is part of the rationality assumptions made precise in
Chapters~\ref{ch:vertex_algebras} and~\ref{ch:regularity}. Suppose
for simplicity that the left and right chiral algebras agree, with
irreducible modules $M_i$. The full state space then decomposes as
\[
 \mathscr H=\bigoplus_{i,j}\C^{Z_{ij}}\otimes
                          M_i\otimes\overline{M_j}.
\]
This formula describes finite-energy states; the physical Hilbert
space includes their Hilbert completion. The barred module denotes
the corresponding antiholomorphic factor. The multiplicity $Z_{ij}$
counts how many copies of the left--right pair occur and is therefore
a nonnegative integer. Taking a trace on each tensor factor gives
the full torus partition function
\[
 Z(\tau,\bar\tau)=\sum_{i,j}Z_{ij}
      \chi_i(\tau)\overline{\chi_j(\tau)}
\]
where $\chi_i=\chi_{M_i}$. Different left and right algebras require
two different character families, but the same trace argument.
The index $0$ labels the vacuum module. If it is the only sector
with weight zero, a unique vacuum requires $Z_{00}=1$.
The simple choice $Z_{ij}=\delta_{ij}$ pairs each left sector with
its corresponding right sector and gives
$Z=\sum_i|\chi_i|^2$; it is the diagonal pairing in this convention.

Changing the oriented basis of the torus's two fundamental cycles
generates $\tau\mapsto\tau+1$ and $\tau\mapsto-1/\tau$.
Neither cycle is contractible on the torus. The first transformation
changes one basis cycle by adding the other; the second exchanges
the cycles with an orientation sign. They generate the modular
group acting on $\tau$. A bosonic full theory with the anomaly
cancelled between left and right must assign the same partition
function to these equivalent descriptions of the torus. This
\emph{modular invariance} constrains the multiplicity matrix
$(Z_{ij})$, but it does not give the three-point constants.
Different local theories can have the same torus state count.

\subsection{The thermal meaning of the trace}

For a thermal circle of Euclidean-time circumference $\beta$, the
inverse temperature in units with Boltzmann's constant one, take
$\tau=iv\beta/L$. Then
$Z=\Tr_{\mathscr H}e^{-\beta H_{\rm cyl}}$.
Indeed multiplying $q^{L_0-c/24}$ by its right-moving counterpart
gives the Boltzmann factor from \eqref{eq:cft-cylinder}. More generally,
writing $\tau=\tau_1+i\tau_2$, the exponent in that full trace is
\[
 -2\pi\tau_2\left(L_0+\bar L_0-\frac{c+\bar c}{24}\right)
 +2\pi i\tau_1\left(L_0-\bar L_0-\frac{c-\bar c}{24}\right).
\]
Thus $\tau_2$ weights energy and $\tau_1$ inserts a spatial
translation, or equivalently a phase determined by momentum.

The modular transformation $\tau\mapsto-1/\tau$ exchanges a short
thermal cycle with a long one. In a unitary \emph{compact} theory,
meaning here a discrete finite-degeneracy spectrum and a convergent
thermal trace, a unique normalizable vacuum controls the
low-temperature limit. It therefore also determines the leading
high-temperature growth through modular invariance. For equal
chiral charges the result is
$\log Z\sim\pi cL/(6v\beta)$ as $v\beta/L\to0$; the exercise
at the end of the chapter derives it from the vacuum term.
In a nonunitary chiral theory whose smallest
weight is $h_{\min}$, the analogous exponent involves the
\emph{effective central charge} $c_{\rm eff}=c-24h_{\min}$,
provided the modularly transformed trace has a nonzero leading
contribution from that lowest-weight sector. For a continuum of states without a
normalizable vacuum, one must analyze the spectral measure instead
of inserting a vacuum term by hand. These qualifications are central
to the later reduced-theory proposals.

\section{Crossing and the meaning of a conformal block}
\label{sec:cft-blocks}
\index{crossing symmetry}\index{conformal block}

\subsection{A channel is a choice of intermediate states}

Four distinct insertion points contain one complex variable that
cannot be removed by a global conformal transformation. This
\emph{cross-ratio} can be written
\[
 x=\frac{z_{12}z_{34}}{z_{13}z_{24}}.
\]
A M\"obius transformation $z\mapsto(az+b)/(cz+d)$, with
$ad-bc\ne0$, puts the ordered points at $(\infty,1,x,0)$.
Here $a,b,c,d$ are coordinate-map coefficients, unrelated to the
central charge. We use the chiral correlator
\[
 G(x)=\langle\phi_1(\infty)\phi_2(1)\phi_3(x)\phi_4(0)\rangle
\]
with the infinity normalization already defined. Near $x=0$ we
expand the pair at $x,0$; near $x=1$ we expand the pair at $1,x$.
These expansions, analytically continued with consistent branches,
must describe the same correlator. This is \emph{crossing symmetry}.
For full fields the corresponding function is $G(x,\bar x)$.

A \emph{conformal block} sums the intermediate states in one
representation of the chosen symmetry algebra, with primary
couplings factored out. To see what is being summed, choose at level
$N$ a basis $|p,N,a\rangle$ of descendants of an exchanged
weight-$p$ state, after removing null vectors. If $G_N$ is its Gram
matrix, the projector onto that level is
\[
 \Pi_{p,N}=\sum_{a,b}|p,N,a\rangle
                    (G_N^{-1})_{ab}\langle p,N,b|.
\]
Insert this projector between the pair of fields inside the radial
circle and the pair outside. Each term is a product of two
three-point matrix elements contracted with an inverse norm matrix.
Summing over $N$ gives the block. This explains why descendant
norms, null quotients, and OPE coefficients enter the same calculation.

\subsection{Global blocks and Virasoro blocks}

The global algebra is generated by $L_{-1},L_0,L_1$. Its
highest-weight state need only be quasiprimary with respect to
Virasoro, and its descendants are $L_{-1}^n|p\rangle$.
For identical chiral external weights $h$ and a unit-normalized
exchanged state of weight $p>0$, the squared norm at level $n$ is
$n!(2p)_n$. After the primary couplings are factored
out, each three-point matrix element contributes $(p)_n$.
Consequently the global block in our unreduced correlator is
\[
 x^{p-2h}\,{}_2F_1(p,p;2p;x).
\]
Here the series is explicitly
\[
 {}_2F_1(p,p;2p;x)
 =\sum_{n\ge0}\frac{(p)_n^2}{(2p)_n\,n!}x^n,
 \qquad (p)_n=p(p+1)\cdots(p+n-1),\quad(p)_0=1.
\]
It converges for $|x|<1$ and sums $L_{-1}$ descendants;
continuation to other channels requires a choice of path.
Multiplying $G(x)$ by $x^{2h}$ defines a reduced correlator. Its
corresponding block is $x^p\,{}_2F_1(p,p;2p;x)$; the different
prefactor is only a convention. The global identity block is
$x^{-2h}$ in the unreduced convention, since the identity has no
nonzero $L_{-1}$ descendants. One should not substitute $p=0$
into the displayed fractions without treating this case separately.
A Virasoro block
also sums descendants involving $L_{-2},L_{-3},\ldots$, and depends
on $c$. A block for an extended algebra sums still more states.
These three expansions are related but are not interchangeable.
In particular, the Virasoro vacuum block includes the stress state
at level two, even though the global identity block has no such
descendant.

\begin{example}[An elementary global block]
\label{ex:cft-global-logarithm}
At exchanged weight $p=1$, $(1)_n=n!$ and $(2)_n=(n+1)!$.
Thus the series coefficient is $1/(n+1)$ and
\[
 x^{1-2h}\sum_{n\ge0}\frac{x^n}{n+1}
 =-x^{-2h}\log(1-x)
 =x^{1-2h}\left(1+\frac{x}{2}+\frac{x^2}{3}+\cdots\right).
\]
The logarithm here is the ordinary analytic function on the branch
reached from $x=0$. Its branch point at $x=1$ shows why continuation
between channels requires more than merely substituting a new real
number into the original power series. This is a global block, not
the complete four-point function or its Virasoro completion.
\end{example}

\emph{Reflection positivity} means that reflecting an inner set of
insertions through the quantization circle and taking their adjoints
produces an outer set whose correlator with the inner set is a
nonnegative squared norm. With several choices of inner state it
produces a positive semidefinite Gram matrix.
In such a four-point channel, exchanged primaries
carry nonnegative coefficients such as $|C|^2$ after orthonormalizing
multiplicity spaces. Together with crossing, this positivity gives
inequalities constraining allowed weights and OPE coefficients, a
method called the \emph{conformal bootstrap}. In a mixed charged channel, individual coefficients can
instead be complex products of distinct OPE constants. Positivity
must be applied to the correct matrix of correlators, not imposed
on every coefficient separately.

When a rational chiral theory has the required tensor-category
structure, changing the chiral fusion channel is encoded by an
associator, of the kind studied in Chapter~\ref{ch:fusion}. Its
finite-dimensional spaces label independent fusion couplings; each
block separately sums an infinite family of descendant states.
For the Haagerup construction the category of chiral modules and
the category of symmetry lines are distinct categories related by
the construction. Their associator matrices cannot simply be
identified. Chapter~\ref{ch:sewing} develops that distinction.

Crossing is a local consistency test, while modular invariance is a
global torus consistency test. A construction must eventually satisfy
both, as well as the other \emph{sewing identities}: compatibility
of correlators when surfaces are cut into simpler pieces and glued
back together in different orders.
Starting with an actual lattice vertex algebra supplies chiral
associativity: the different orders of expanding its OPEs agree.
It does not remove the work of identifying its representation
category or the full two-dimensional theory built from it.

\section{Perturbations and what dimension can tell us}
\label{sec:cft-perturbations}
\index{relevant perturbation}\index{marginality}

Let $S_*$ be the Euclidean action of a CFT. A \emph{perturbation}
changes it to
\[
 S=S_*+\lambda_{\rm phys}\int d^2x\,\Phi(x),
\]
where $\Phi$ is a scalar of dimension $\Delta$. The action is
dimensionless, the area element has units of length squared, and
$\Phi$ has units of length to the power $-\Delta$.
Consequently $\lambda_{\rm phys}$ has units of
length$^{\Delta-2}$, or mass dimension $2-\Delta$.
At an observation length $\ell$, its dimensionless strength at
leading order is $g(\ell)=\lambda_{\rm phys}\ell^{2-\Delta}$.

The \emph{renormalization group} tracks how the effective
dimensionless couplings change when short-distance fluctuations are
averaged out and the observation scale is increased. The derivative
of a running coupling with respect to $\log\ell$ is its
\emph{beta function}. Including interaction corrections, and using
the convention that $\ell$ increases toward long distances, it has
the form
\[
 \frac{dg}{d\log\ell}
 =(2-\Delta)g+b g^2+\cdots.
\]
Thus $\Delta<2$ is \emph{relevant}: a small perturbation grows
at long distances to first order. For $\Delta>2$ it decreases and
is \emph{irrelevant}. When $\Delta=2$ it is \emph{marginal by
dimension}; the linear term gives no answer about its eventual flow.
The quadratic coefficient depends on normalized OPE data and on the
coupling convention. For a single marginal coupling, keeping only
that term and taking $g(\ell_0)=g_0$ gives
\[
 g(\ell)=\frac{g_0}{1-bg_0\log(\ell/\ell_0)}.
\]
This approximation is valid while $|g|$ remains small. If $bg_0<0$
the coupling approaches zero logarithmically (marginally irrelevant);
if $bg_0>0$ its magnitude grows (marginally relevant). If $b=0$,
higher terms can still produce running. An \emph{exactly marginal}
deformation has vanishing beta functions to all orders and gives
a continuous family of CFTs. Several dimension-two fields can mix, and
symmetry can permit or forbid their coefficients independently of
their dimensions.
Worked examples~\ref{ex:full-nonabelian-running} and
\ref{ex:full-current-logarithms} derive a nonzero quadratic term
from current OPEs and show how it produces slow logarithmic drift.

The \emph{scalar gap} is the smallest positive scaling dimension
of a nonidentity spin-zero primary in the specified spectrum. It is
a dimension gap, not a nonzero energy gap in the infinite-volume
Hamiltonian. The \emph{symmetry-invariant scalar gap} restricts the
same question to symmetry-preserving perturbations. For a group this
means that every group element leaves the field unchanged. For a
noninvertible line $a$, invariance means that its line operator acts
on the field's state by the vacuum eigenvalue $d_a$, the quantum
dimension, as explained in Chapter~\ref{ch:fusion}.
A relevant
scalar may exist but be forbidden in a symmetry-preserving
Hamiltonian. Conversely, a gap of two in the invariant sector does
not by itself prove a stable \emph{conformal manifold}, a continuous
family of inequivalent conformal theories parametrized by exactly
marginal couplings. One must solve the coupled beta functions and
identify deformations that merely relabel the same theory. We will compute
these different gaps explicitly for the two full $c=8$ theories.

\section{Exercises}

\begin{exercise}\label{ex:cft-level1}
Use the Virasoro commutator to show
$\norm{L_{-1}^n|h\rangle}^2=n!(2h)_n$ for a unit primary.
Here $(a)_n=a(a+1)\cdots(a+n-1)$. Derive the first two coefficients
of the global block from these norms.

\par\smallskip\noindent{\footnotesize\normalfont
\hyperref[sol:ex:cft-level1]{Solution on page~\pageref*{sol:ex:cft-level1}.}}\par
\end{exercise}

\begin{exercise}\label{ex:cft-null}
At $c=1/2,h=1/16$, find the null vector at level two from
\eqref{eq:cft-gram2}. Check directly that $L_1$ and $L_2$ annihilate it.

\par\smallskip\noindent{\footnotesize\normalfont
\hyperref[sol:ex:cft-null]{Solution on page~\pageref*{sol:ex:cft-null}.}}\par
\end{exercise}

\begin{exercise}\label{ex:cft-compact}
In the compact-boson convention of this chapter, determine the radius
at which the $(m,n)=(1,1)$ field is a holomorphic current. Show that
the transformation $R\mapsto2/R$, $m\leftrightarrow n$ preserves
the spectrum.

\par\smallskip\noindent{\footnotesize\normalfont
\hyperref[sol:ex:cft-compact]{Solution on page~\pageref*{sol:ex:cft-compact}.}}\par
\end{exercise}

\begin{exercise}\label{ex:cft-stress-combination}
Let $T_1,T_2$ be commuting stresses of positive central charges
$c_1,c_2$, meaning that their mixed OPE has no singular part.
Prove that $A=c_2T_1-c_1T_2$ is a weight-two primary
for $T=T_1+T_2$, and compute its norm. Explain why $A$ is not
itself a stress tensor with this normalization.

\par\smallskip\noindent{\footnotesize\normalfont
\hyperref[sol:ex:cft-stress-combination]{Solution on page~\pageref*{sol:ex:cft-stress-combination}.}}\par
\end{exercise}

\begin{exercise}\label{ex:cft-thermal}
For a unitary compact theory with $c=\bar c$, a discrete spectrum,
a convergent thermal trace, and a unique normalizable vacuum,
derive the leading high-temperature
formula $\log Z\sim\pi cL/(6v\beta)$ from modular invariance and
vacuum dominance at low temperature. Identify the step that fails when the
vacuum is absent from the normalizable Hilbert space.

\par\smallskip\noindent{\footnotesize\normalfont
\hyperref[sol:ex:cft-thermal]{Solution on page~\pageref*{sol:ex:cft-thermal}.}}\par
\end{exercise}

\begin{exercise}\label{ex:cft-marginality}
A character calculation finds thirty-six scalar primaries of
dimension two invariant under a specified symmetry. State three
additional kinds of information needed before asserting that they
parametrize a thirty-six-dimensional conformal manifold.

\par\smallskip\noindent{\footnotesize\normalfont
\hyperref[sol:ex:cft-marginality]{Solution on page~\pageref*{sol:ex:cft-marginality}.}}\par
\end{exercise}
